%% kapa: exam=1 round=23 year=2012 period=1 form=A subject=civil_law subject_ko=민법 seq=12 number=12
%% source: Q-Net 감정평가사 기출문제 (cst00309) artlSeq=5052578 file=raw/1008/5052578/2012년 제23회 감정평가사 제1차시험 기출문제(1교시 A형).pdf
%% conversion: kordoc 4.24.51 -> tools/gen1.py (auto); review=proofread-sonnet; flags=none
\begin{question}{12}
신의성실의 원칙과 권리남용 금지원칙에 관한 설명으로 옳지 않은 것은? (다툼이 있으면 판례에 의함)
\begin{choices}
  \item 강행법규 위반임을 알면서 계약을 체결한 후 이를 이유로 계약의 무효를 주장하는 것은 특별한 사정이 없는 한 신의칙에 반한다.
  \item 실효의 원칙을 적용하기 위해서는 의무자인 상대방이 더 이상 권리자가 그 권리를 행사하지 아니할 것으로 믿을 만한 정당한 사유가 있어야 한다.
  \item 채무자가 채권자의 시효중단을 현저히 곤란하게 하여 채권자가 시효중단을 주장할 수 없는 경우, 채무자가 시효의 완성을 주장하는 것은 신의칙상 허용되지 않는다.
  \item 법정지상권을 가진 건물소유자로부터 건물을 양수하면서 법정지상권까지 양도받기로 한 자에 대하여 대지소유자가 그 소유권에 기하여 건물철거를 청구하는 것은 허용되지 않는다.
  \item 권리남용의 해당 여부는 당사자의 주장이 없더라도 법원이 직권으로 판단할 수 있다.
\end{choices}
\end{question}
\answer{1}
